\documentclass[11pt,a4paper]{article}

\usepackage[margin=1in]{geometry}
\usepackage[T1]{fontenc}
\usepackage[utf8]{inputenc}
\usepackage{lmodern}
\usepackage{microtype}
\usepackage[round,authoryear]{natbib}
\usepackage[colorlinks=true,linkcolor=blue,citecolor=blue,urlcolor=blue]{hyperref}

\title{Using ChatGPT Voice for spoken German practice:\\A case study}
\author{%
  \parbox{0.96\textwidth}{\centering
    Sarah Wright\textsuperscript{a}, Manny Rayner\textsuperscript{b},
    ChatGPT C-LARA-Instance\textsuperscript{c}, and Branislav B\'{e}di\textsuperscript{d}\\[0.6em]
    \small \textsuperscript{a}Independent researcher, Australia,
    \href{mailto:sarah@tamargas.com.au}{sarah@tamargas.com.au}\\
    \small \textsuperscript{b}Independent researcher, Australia,
    \href{mailto:mannyrayner@c-lara.org}{mannyrayner@c-lara.org}\\
    \small \textsuperscript{c}OpenAI large language model instance participating in the C-LARA project,
    \href{mailto:chatgptclarainstance@proton.me}{chatgptclarainstance@proton.me}\\
    \small \textsuperscript{d}The \'{A}rni Magn\'{u}sson Institute for Icelandic Studies, Iceland,
    \href{mailto:branislav.bedi@arnastofnun.is}{branislav.bedi@arnastofnun.is}
  }%
}
\date{2026}

\begin{document}

\maketitle

\begin{abstract}
This article reports a case study in which two anglophone learners used OpenAI's ChatGPT Pro Tier Voice Mode for German conversation practice approximately once a week over eight weeks, usually in 60--90 minute sessions. The sessions used a loosely structured microlearning environment, initially organised around practical travel role-play scenarios, and gradually developing into broader and more personally engaging conversations. Participant observations and post-session summaries suggest that the system supported confidence, fluency, situated vocabulary learning, flexible topic development, and low-pressure practice. On the negative side, important limitations were apparent: the system lacked corrective feedback, consistent identification of speakers, effective clarification strategies, systematic learner tracking, and fact checking. The most important practical consideration at the moment is the high subscription cost, which presents a substantial access barrier. Although our study is small-scale and only presents anecdotal evidence, it documents a simple, reproducible use of a generally available commercial AI system and identifies priorities for further work which could potentially lead to transformative changes in spoken CALL in the short to medium term.
\end{abstract}

\noindent\textbf{Keywords:} spoken interaction; ChatGPT Voice Mode; German language learning; foreign-language anxiety; conversational AI.

\section{Introduction and overview}

Learners often have only limited opportunities for spoken language practice in foreign and second language learning (L2), both in classroom environments and during self-study. Even learners who have studied an L2 for years and accumulated substantial passive knowledge may still, as a result of inadequate spoken practice, lack the confidence to use the language spontaneously. Many modern technologies support personalised and adaptive L2 learner experiences while providing opportunities for real-time interaction in both written and spoken form \citep{Urbaite2024}. In particular, technologies powered by generative artificial intelligence (GenAI) not only assist learners in meeting their overall learning goals in L2 but are also reported as contributing to increased speaking proficiency and reduced speaking anxiety \citep{DuDaniel2024,MengZou2026,SaarelaEtAl2026}.

Our article focuses on an experiment involving the use of OpenAI ChatGPT Voice Mode as a conversation partner to support L2 German speaking skills in unconstrained situations. The Pro Tier version we used is currently priced at USD~200 per month, and is primarily designed for coders and business users rather than as an educational language-learning product \citep{Reuters2024}. Although this cost represents a practical barrier to adoption, and could at present be regarded as contributing to the digital divide in language learning \citep{Yaman2015}, the platform is publicly available, which makes our study straightforward to replicate.

We should note two things at once. First, the case study is entirely anecdotal. It involves only one main subject, who needed to improve her spoken German conversational skills for practical reasons. She knew that one of the authors had access to ChatGPT Voice Mode / Pro Tier, and asked whether it would be possible to use it for these purposes; based on the remarkably positive results we obtained even in the first session, we decided to keep records so that we could write up the outcomes. Second, it became clear that Pro Tier appears in practice to be necessary for the type of free conversation practice we describe here. We experimented first with the much cheaper Plus Tier (USD~20 per month), and found that performance was clearly inadequate. As we explain in the next section, one of the reasons why the experiment seemed worth writing up is that we have not found any clearly comparable studies reported in the CALL literature.

The work in our experiment is related to the broader C-LARA project, a ChatGPT-based online open-source language-learning platform supporting the creation of multimedia language-learning resources, including glossed texts, illustrated picture books, audio, and learner-specific resources \citep{BediEtAl2025}. Another paper we present at this conference describes C-LARA-2, a reimplemented version with rationalised and improved architecture. We would ideally like to include spoken interaction in C-LARA-2; the work we describe here should thus be viewed not only as contributing to a general understanding of how functionalities like ChatGPT Voice Mode may enhance spoken interaction skills, but also as a pre-study to investigate whether AI-based spoken dialogue functionalities of this kind could possibly become part of C-LARA-2.

Methodologically, our case study is situated within the framework of individualised and self-regulated learning, which helps learners adapt to individual differences, manage their own learning processes, and gradually become more independent while achieving their learning goals \citep{Zimmerman2002}. Individual sessions using the ChatGPT Voice Mode app create a unique microlearning environment for enhancing specific language skills \citep{KohnkeMoorhouse2026}. Building on this view of self-regulated learning, the study examines whether ChatGPT Voice Mode can support speaking and listening practice during home self-study and partly take over the role of a human conversation tutor. This is explored through six research questions: (1) whether ChatGPT Voice Mode is sufficiently good at supporting realistic conversations in German; (2) whether it can adjust to the learners' level and maintain a conversation while providing input that is not too difficult for them to understand; (3) whether it can recognise what the learners say and respond accordingly; (4) how well the Voice Mode feature can maintain the focus of the conversation by asking follow-up questions, encouraging turn-taking, rephrasing when the learner does not understand, and keeping the conversation interesting on a topic selected by the learner; (5) whether regular use of ChatGPT Voice Mode can improve learners' confidence and fluency in spoken German; and (6) how a Voice Mode feature could improve existing CALL online platforms such as C-LARA.

The rest of the paper is organised as follows. Section~2 describes previous research in GenAI-supported spoken interaction in CALL. Section~3 describes the background to the study. Section~4 presents findings. Section~5 discusses how Voice Mode could be integrated with text-based CALL functionalities, in particular with our own C-LARA-2 platform. The final section summarises and concludes.

\section{Previous work}

Several recent papers suggest that combining spoken interaction and GenAI in CALL is a promising development that is starting to attract considerable interest. For example, \citet[p.~122]{Savvidou2026} claims that GenAI-powered CALL platforms like ChatGPT Voice Mode, Duolingo Max, Gemini, Speak AI, Replica AI, and Talkpal AI can be helpful tools offering realistic speaking-practice opportunities, immediate feedback, personalised conversations, and opportunities for learners to continue practising independently outside the classroom, but unfortunately gives few details. In a study conducted by \citet{LuEtAl2026}, ChatGPT generates conversational scenarios while spoken responses from English language learners are captured by a speech-recognition system and the conversation is conducted with an animated conversational agent generated by another system called D-ID AI. Their results claim that repeated practice improves learners' speaking proficiency in the target language and reduces anxiety from having real-life conversations, but it is unclear from their paper what kind of spoken interaction is actually supported. The study we know of which is closest to ours is that of \citet{CarreraNunezEtAl2025}, which describes an experiment in which ChatGPT Voice Mode was incorporated into L2 learning in regular and guided practice sessions in a classroom environment. Tests provided strong statistical evidence that it led to improvement of different language skills, particularly pronunciation, fluency, vocabulary, and grammar, as well as the reduction of language anxiety; the authors attribute this to the tool's non-judgemental environment.

All of the studies cited above stress how powerful the combination of GenAI and spoken interaction can be, particularly for learners with an autonomous learning style, and their overall conclusions agree well in principle with the findings we present here. The details, however, are less clear to us: crucially, we are uncertain about the extent to which the spoken conversations can reasonably be considered to be unconstrained, an issue which in practice is strongly related to the reliability of the spoken interaction. For obvious reasons, and particularly for maintaining student engagement, this is an important question; as already noted, we obtained clearly unsatisfactory results for unconstrained free conversation with anything other than the highest-level (Pro Tier) version of ChatGPT Voice Mode. The technology concerned is evolving so quickly that this may only reflect a transient issue, but we note it for what it is worth.

\section{Background to the study}

As we said in the previous section, there is rapidly growing interest in the use of generative AI tools for speaking practice outside the classroom, but it is unclear to us that the most commonly available platforms can adequately support unconstrained conversation and function as true conversation partners during home self-study. Since our preliminary experiments suggested that this was not the case, the central research goal of the study reported here was to explore whether a high-end system, ChatGPT Pro Voice Mode, would be equal to the task. In particular, we wished to know whether it would allow learners to take part in spoken interaction, practise everyday situations, ask for vocabulary, and receive immediate responses.

The experiments used a small-scale qualitative case-study design and were carried out during eight regular sessions over approximately eight weeks, each lasting about 60--90 minutes. Participant~1, the main learner, was in her early twenties. She was a native speaker of English who had previously studied German for six years and been awarded a B1-level certificate reflecting reasonable low-intermediate skills in reading, writing, speaking, and listening; unfortunately, she had not maintained contact with the language for some time, and in particular felt that her fluency had declined considerably. She was, however, planning a trip where she would spend an extended period in German-speaking countries, and consequently wanted to improve her confidence in spoken German. For this reason, we agreed that we would begin by practising conversations organised around practical scenarios she could expect to encounter on her trip, and, if these went well, progress to more challenging topics.

Participant~2, in his mid-sixties, was a native speaker of English who had previously lived in several countries and worked in numerous jobs involving speech and language technology. He had near-native skills in Swedish, strong (C1) reading ability in French, Norwegian, and Danish, and B1/B2 reading ability in German, but his spoken German skills were only around A2 level. He participated as a facilitator and observer, helping to guide the conversation when needed; his main motivation for improving his German was to be able to watch German films more easily. The third participant, who acted as the external observer, was an experienced language teacher. He attended one early session and followed the remaining sessions through written summaries.

The sessions were not recorded or formally transcribed. After each session, ChatGPT generated a detailed written summary, which formed the main dataset. These summaries were analysed using directed qualitative content analysis \citep{AssarroudiEtAl2018}, which helped to extract relevant data based on predefined research criteria: the naturalness and real-life quality of the conversation, adaptation to the learner's language level, accuracy of speech recognition and appropriateness of responses, conversation management and learner engagement, and possible changes in the learners' confidence and fluency.

\section{Findings}

The sessions suggested several ways in which ChatGPT Voice Mode / Pro Tier can support spoken language practice. Since the study was only intended to be exploratory and offered anecdotal evidence, these benefits could not be measured quantitatively. They were, however, visible as the sessions unfolded and as the participants reflected on their learning experience. This section provides an interpretation of the findings divided into two categories: observed benefits and limitations.

\subsection{Observed benefits}

\subsubsection*{Authentic conversation quality leading to increased learner confidence}

The most striking benefit was Participant~1's increased willingness to speak. At the beginning of the experiment, her difficulty was not primarily a lack of German vocabulary or grammar. She had studied German at school and retained a reasonable passive understanding, but was hesitant about using the language spontaneously. ChatGPT Voice Mode / Pro Tier provided a setting in which she could speak without the social pressure of addressing an unfamiliar human interlocutor. It generally treated partially successful utterances as meaningful and responded to the intended content, thus minimising the stress associated with the fear of speaking incorrectly and, as a consequence, creating a low-pressure environment for fluency practice. After only two or three weeks, Participant~1 was already less dependent on rehearsed scenarios and more willing to initiate topics, ask questions, and continue speaking despite uncertainty.

The sessions provided a substantial amount of practice. Meeting approximately weekly for 60--90 minutes gave Participant~1 extended exposure to real-time spoken interaction. This is difficult to obtain in many ordinary learning contexts. A human conversation partner may not be available, and a teacher's time is usually limited. ChatGPT Voice Mode / Pro Tier could sustain long conversations without a fixed lesson plan and without the fatigue or scheduling constraints that often restrict spoken practice.

\subsubsection*{ChatGPT Voice Mode adjusting to the learner's level by providing scaffolding}

A second important benefit was immediate scaffolding. When Participant~1 needed a word or phrase, ChatGPT Voice Mode / Pro Tier could usually supply it quickly and naturally within the conversation. Examples included practical vocabulary such as \emph{der Eiffelturm} for ``the Eiffel Tower'', \emph{das Armband} for ``bracelet'', \emph{die Aussprache} for ``pronunciation'', \emph{der Dieb} and \emph{der Diebstahl} for ``thief'' and ``theft'', \emph{das Katzenfutter} and \emph{das Katzenleckerli} for ``cat food'' and ``cat treats'', \emph{der Thunfisch} for ``tuna'', and \emph{die \'{U}berschwemmung} for ``flooding''. These items were not introduced as part of a pre-planned syllabus. They arose, rather, because the learner needed them in order to say something she spontaneously wanted to express. Voice Mode also helped the learner remain in the target-language interaction when she lacked a precise form. In some cases the AI tool supplied a missing word; in others, the tool inferred the intended meaning and responded in simple German. This allowed the conversation to continue without repeatedly stopping for translation. The scaffolding was therefore not only lexical but also interactional.

\subsubsection*{Conversation management and learner engagement supporting role-play scenarios}

A third benefit was the system's flexibility. The early sessions rehearsed practical travel scenarios of the kind informally agreed at the beginning of the study, organised as role plays between the human and the AI (checking into a campsite, visiting a pharmacy, and so on). Later sessions became less scripted and more personal, talking about previous trips, food and drink, music, concerts, chess, fairy tales, films, pets, weather, pronunciation, and AI. This strongly suggests that ChatGPT Voice Mode / Pro Tier was not merely useful for practising set travel phrases, but could from the start support natural conversation. This flexibility is one of the distinctive strengths of a general conversational AI. A conventional coursebook can provide carefully graded scenarios, but it cannot easily follow the learner from a campsite role-play to a discussion of Gershwin, bad weather, cat treats, or \emph{2001: A Space Odyssey}. A human teacher can do this, but has limited time. Voice Mode offered a form of abundant, responsive conversational practice that could adapt to the learner's interests in real time.

\subsection{Observed limitations}

\subsubsection*{Limited pedagogical control}

The most important limitation was the lack of explicit pedagogical control. Voice Mode behaved primarily as a helpful conversation partner. This made the sessions relaxed and encouraging, but it also meant that the system did not consistently pursue language-learning goals. It did not maintain an explicit learner model, track Participant~1's recurring errors, or adapt its correction strategy over time. In particular, the AI usually prioritised communicative success over formal accuracy.

\subsubsection*{Speech recognition, clarification, and turn-taking}

A second set of limitations concerned the mechanics of spoken interaction. The speech input was sometimes noisy, incomplete, or incorrectly interpreted, especially when the speakers hesitated, code-switched, used proper names, or spoke in fragments. In several cases the AI responded to what it believed it had heard rather than to what the speaker had intended. Some misunderstandings were minor, but others temporarily diverted the conversation. Turn-taking also caused occasional problems. The system sometimes treated a pause as the end of a turn, produced a generic response where a more targeted follow-up would have been preferable, or decided to close the conversation prematurely with farewell formulas.

\subsubsection*{Confusing speaker voices}

The sessions usually involved Participant~1 and Participant~2 sitting together and both speaking to ChatGPT. This created a further problem: the system did not reliably know who was speaking, which was sometimes an annoying limitation.

\subsubsection*{Factual reliability}

The conversations often moved beyond language practice into cultural, historical, linguistic, or technical information. In many cases the AI's responses were useful and plausible, but they were not systematically checked during the sessions. This creates a familiar risk: fluent conversational delivery can make uncertain information about diverse topics provided by ChatGPT Voice Mode / Pro Tier sound authoritative. In the opinion of the two participants, however, there were very few occasions where information provided by the AI was in fact dubious or incorrect.

\subsubsection*{Access and cost}

As noted, an extremely important practical limitation right now is cost, since acceptable performance for extended, flexible spoken interaction depends on having access to the expensive ChatGPT Pro Tier. We only had a Pro Tier subscription because we use it for coding in the C-LARA project, and we think it unlikely that many people would purchase one just to use it as a conversation partner. In practice, we suspect that access will be limited to people who, like us, already have it for other reasons.

\section{Further directions: integration of Voice Mode and multimodal texts}

The experiments we have described suggest that it could potentially be very useful to add a spoken-language component to our own C-LARA platform, which supports the creation and use of multimodal pedagogical texts. We have already carried out a preliminary integration experiment. All that is needed is to open two browser windows, one for C-LARA and one for ChatGPT Pro. The user accesses a C-LARA text page, takes a screenshot, and uploads it to the ChatGPT Pro window. They then switch to Voice Mode. It is indeed possible to discuss the content of the C-LARA page, with Voice Mode clearly understanding both the text and the images. This integration is, needless to say, unbearably cumbersome for practical use, but smoother integrations can be envisaged.

A more interesting issue is that the combined system would need to add stronger pedagogical structure. A dedicated language-practice mode could allow learners or teachers to choose how often corrections are given, while delayed feedback could support fluency during the session and accuracy afterwards. The experiment also shows the value of post-session processing. The combined system could automatically turn each conversation into a summary, vocabulary list, corrected sentences, grammar explanations, and follow-up exercises. Topics discussed during the session could then generate personalised dialogues, short texts, glossed passages, audio activities, or reading exercises. This would create a continuous learning cycle in which spoken practice produces new study materials.

Other important areas are learner tracking and teacher guidance. The combined system should distinguish between speakers, connect errors and vocabulary needs to the correct learner, and record progress across several sessions. Conversations could also be linked to visible texts, images, or exercises, giving the interaction a clearer focus. It could also be desirable to have teachers continue to control learning goals, assessment, and final pedagogical decisions.

The idea of creating this kind of integration of C-LARA and Voice Mode is extremely tempting, and it is mainly a question of whether we have the necessary resources. Irrespective of whether we do so or not, the idea is so obviously natural and useful that we would be surprised if someone did not build a platform of this general type soon. For all we know, it may exist already.

\section{Summary and conclusion}

We have presented a small study in which we explored ChatGPT Voice Mode / Pro Tier's ability to support German speaking practice during independent home study. We do not claim that the results are more than anecdotal. They refer to one main subject and are based on AI-generated post-session summaries and informal subject feedback rather than full recordings, transcripts, or independent speaking assessments, and cannot, therefore, demonstrate measurable improvements.

With the above caveats, we are much encouraged by our experience. ChatGPT Voice Mode / Pro Tier showed strong potential for low-pressure conversation practice. It was able to maintain extended and natural conversations about both constrained and open-ended topics. It could adapt its language to the learner's level by simplifying, paraphrasing, explaining, providing vocabulary, and generally reformulating when difficulties occurred. Speech recognition was mostly sufficient for the conversation to continue, although misunderstandings and requests for repetition still appeared, particularly when pronunciation or wording was unclear.

ChatGPT Voice Mode / Pro Tier also showed an ability to manage the conversation by asking follow-up questions, supporting turn-taking, maintaining selected topics, and elaborating its explanation when the learner did not understand. Its flexible role as conversation partner, vocabulary assistant, and provider of clarifications helped the learner to continue speaking without frequent interruption. Over the eight sessions, the learner became more willing to introduce topics, ask questions directly, and continue speaking despite limited vocabulary or grammatical uncertainty. She reported this as experiencing increased confidence and greater conversational fluency, and the observer concurred in this judgement. This is in line with \citet{CarreraNunezEtAl2025}, who find that regular sessions with GenAI Voice Mode over a longer period of time can contribute to increased language skills.

Although our contribution is modest, we interpret it as suggesting that ChatGPT Voice Mode / Pro Tier and similar high-end systems should be taken seriously as existing, already available providers of useful spoken language practice. They not only give instant access to a virtual conversation partner, imitate real-life scenarios, supply vocabulary on request, tolerate imperfect learner language, and allow discussion of personally meaningful topics, but also help reduce speaking anxiety; all of this agrees well with previous research \citep{DuDaniel2024,MengZou2026,SaarelaEtAl2026}. Since the platform we used is openly available, anyone can repeat our experiments. We would very much like to see other people doing so and reporting more substantial results. Most importantly, we wish to know whether they concur in our judgement that it can support long, unconstrained, pedagogically useful conversations with a non-trivial and varied sample of learners.

As we have said, although the mid-2026 version of ChatGPT Voice Mode / Pro Tier is already very useful as a conversation partner, it has clear limitations. It is not a specialised CALL tutor. It prioritises conversational flow over systematic correction to an extent that can be counterproductive, especially since it can misinterpret spoken input. It does not reliably distinguish speakers. It is not well integrated with text and images. Above all, it is too expensive for most language learners to consider using it. We think it possible, however, that the picture might soon change, in which case the implications could be considerable. Watch this space.

\section*{Division of work statement}

The role of Participant~1 (main learner) was occupied by Sarah Wright. The role of Participant~2 (co-learner) was occupied by Manny Rayner. The role of the External Observer was occupied by Branislav B\'{e}di. The role of the AI was occupied by ChatGPT C-LARA-Instance.

The first draft of the paper was written by ChatGPT C-LARA-Instance and then heavily revised, first by Branislav B\'{e}di, then by Manny Rayner, and last by Sarah Wright. The human authors take full responsibility for the content.

\section*{Funding}

The ChatGPT Pro Tier subscription used for these experiments, and also for the development of C-LARA-2, is funded privately by Manny Rayner and Cathy Chua.

\bibliographystyle{plainnat}
\bibliography{references}

\end{document}
